\documentclass[twoside]{article}
\setlength{\oddsidemargin}{0.25 in}
\setlength{\evensidemargin}{-0.25 in}
\setlength{\topmargin}{-0.6 in}
\setlength{\textwidth}{6.5 in}
\setlength{\textheight}{8.5 in}
\setlength{\headsep}{0.75 in}
\setlength{\parindent}{0 in}
\setlength{\parskip}{0.1 in}

%
% ADD PACKAGES here:
%

\usepackage{amsmath}
\usepackage{bbm}
\usepackage{amsfonts}
\usepackage{graphicx}

%
% The following commands set up the lecnum (lecture number)
% counter and make various numbering schemes work relative
% to the lecture number.
%
\newcounter{lecnum}
\renewcommand{\thepage}{\thelecnum-\arabic{page}}
\renewcommand{\thesection}{\thelecnum.\arabic{section}}
\renewcommand{\theequation}{\thelecnum.\arabic{equation}}
\renewcommand{\thefigure}{\thelecnum.\arabic{figure}}
\renewcommand{\thetable}{\thelecnum.\arabic{table}}

%
% The following macro is used to generate the header.
%
\newcommand{\lecture}[4]{
   \pagestyle{myheadings}
   \thispagestyle{plain}
   \newpage
   \setcounter{lecnum}{#1}
   \setcounter{page}{1}
   \noindent
   \begin{center}
   \framebox{
      \vbox{\vspace{2mm}
    \hbox to 6.28in { {\bf SOQ 2026: Summer of Quant
	\hfill May-Jul 2026} }
       \vspace{4mm}
       \hbox to 6.28in { {\Large \hfill Lecture #1: #2  \hfill} }
      \vspace{2mm}}
   }
   \end{center}
   \markboth{Lecture #1: #2}{Lecture #1: #2}
}
%
% Convention for citations is authors' initials followed by the year.
% For example, to cite a paper by Leighton and Maggs you would type
% \cite{LM89}, and to cite a paper by Strassen you would type \cite{S69}.
% (To avoid bibliography problems, for now we redefine the \cite command.)
% Also commands that create a suitable format for the reference list.
\renewcommand{\cite}[1]{[#1]}
\def\beginrefs{\begin{list}%
        {[\arabic{equation}]}{\usecounter{equation}
         \setlength{\leftmargin}{2.0truecm}\setlength{\labelsep}{0.4truecm}%
         \setlength{\labelwidth}{1.6truecm}}}
\def\endrefs{\end{list}}
\def\bibentry#1{\item[\hbox{[#1]}]}

%Use this command for a figure; it puts a figure in wherever you want it.
%usage: \fig{NUMBER}{JUST-THE-FILENAME}{CAPTION}
\newcommand{\fig}[3]{
% 			\vspace{#2}
			\begin{center}
			\includegraphics{figures/#2} 
			\newline
			Figure \thelecnum.#1:~#3
			\end{center}
	}
	
% Use these for theorems, lemmas, proofs, etc.
\newtheorem{theorem}{Theorem}[lecnum]
\newtheorem{lemma}[theorem]{Lemma}
\newtheorem{proposition}[theorem]{Proposition}
\newtheorem{claim}[theorem]{Claim}
\newtheorem{corollary}[theorem]{Corollary}
\newtheorem{definition}[theorem]{Definition}
\newenvironment{proof}{{\bf Proof:}}{\hfill\rule{2mm}{2mm}}

% **** IF YOU WANT TO DEFINE ADDITIONAL MACROS FOR YOURSELF, PUT THEM HERE:

\newcommand\E{\mathbb{E}}

\begin{document}
%FILL IN THE RIGHT INFO.
%\lecture{**LECTURE-NUMBER**}{**DATE**}{**LECTURER**}{**SCRIBE**}
\lecture{1}{Probability, Stochastic Calculus and Basics of Finance}
%\footnotetext{These notes are partially based on those of Nigel Mansell.}


\section{Random Variable} % Don't be this informal in your notes!

A random variable is a number attached to a random outcome.
Each random variable is fully described by its distribution: the values it can take, and how likely it is.

There are two types of random variables:

\begin{enumerate}
   \item[\bf Discrete :] countable values (usually integers). Number of heads in 10 flips, number of trades made today.
   \item[\bf Continuous :] values over a range. Tomorrow's temperature, the price of a stock at noon. 
\end{enumerate}

\subsection{Probability Mass Function}

It lists the probability corresponding to each value.
\begin{align*}
   p(x) &= P(X = x) \\
   \sum_x p(x) &= 1
\end{align*}

\subsection{Probability Density Function}

For a continuous variable, probability of x being a single number is 0. Instead we use probability density $f(x)$ over a domain $D$. Probability of x being between two numbers $a$ and $b$ is:
\begin{align*}
   P(a \leq X \leq b) &= \int_a^bf(x) dx \\
   \int_D f(x) dx &= 1
\end{align*}

\subsection{Cumulative Density Function}

For a continuous variable, the cumulative density function encodes the probability of being at most x.

It is denoted by $F(x)$.
\[
   F(x) = P(X \leq x)
\]
The graph of $F(x)$ vs $x$ is a non-decreasing curve, with $F(x) = 0$ at the extreme left and $F(x) = 1$ at the extreme right.

Also, in the continuous case,
\[
   F'(x) = f(x)
\]

\section{Expectation}

The expectation of a discrete random variable is defined as:
\[
   E[X] = \sum_x x \cdot p(x)
\]

For a continuous variable, the expectation is defined as:
\[
   E[X] = \int_D x \cdot f(x) dx
\]

{\bf Example: } In the case of a fair die, each face has an equal probability of $\frac{1}{6}$.
\[
   E[X] = \frac{1+2+3+4+5+6}{6} = \frac{21}{6} = 3.5
\]

\subsection{Properties of Expectation}

\begin{itemize}
   \item {\bf Linearity}
   
   For any two random variables $X$ and $Y$ and any two constants $a$ and $b$,
   \[
      E[aX+bY] = a \cdot E[X] + b \cdot E[Y]
   \]

   This is true {\bf regardless of whether $X$ and $Y$ are independent or not.}

   \item {\bf The Indicator Trick}
   
   Consider an event $A$. Define another event ${\mathbbm 1_A} = 1$, if $A$ occurs and $0$ otherwise. Then,
   \[
      E[{\mathbbm 1_A}] = 1 \cdot P(A) + 0 \cdot (1 - P(A)) = P(A)
   \]
   Thus, expectation of an indicator event is the probability of the corresponding event.

   If you want the expected number of events that occur out of $A_1, \dots, A_n$, just add up their probabilities. No independence is necessary in this formula.

   \[
      E[\text{Number of $A_i$ that occur}] = P(A_1) + \dots + P(A_n)
   \]

   \item {\bf Composition}
   
   The expected value of a function $g(X)$ of a random variable $X$ is defined as,

   For a discrete random variable:

   \[
      E[g(X)] = \sum_x g(x) \cdot p(x)
   \]

   For a continuous random variable:

   \[
      E[g(X)] = \int_D g(x) \cdot f(x) dx
   \]
\end{itemize}

\section{Variance}

The variance of a random variable is defined as:
\[
   Var(X) = E[(X - \mu)^2]
\]
where,
\[
   \mu = E[X]
\]

Simplifying the expression further:

\begin{align*}
   Var(X) &= E[ (X - \mu)^2 ] \\
   &= E[X^2 - 2 \cdot \mu \cdot X + \mu^2] \\
   &= E[X^2] - 2 \cdot \mu \cdot E[X] + \mu^2 \\
   &= E[X^2] - 2 \cdot \mu^2 + \mu^2 \\ 
   &= E[X^2] - \mu^2 \\ 
   Var(X) &= E[X^2] - (E[X])^2
\end{align*}

The standard deviation $\sigma$ is defined as
\[
   \sigma = \sqrt{Var(X)}
\]

$X$ is typically within $\pm \sigma$ of its expected value.

\subsection{Properties of Variance}

\begin{itemize}
   \item{\bf Transformation and Scaling}
   \begin{align*}
      Var(X + b) &= Var(X) \\
      Var(aX) &= a^2Var(X) \\
      Var(aX+b) &= a^2Var(X)
   \end{align*}

   \item{\bf Variance of sum of two random variables}
   \[
      Var(X + Y) = Var(X) + Var(Y) + 2 \cdot Cov(X, Y)
   \]

   Here, $Cov(X, Y)$ is called the covariance of $X$ and $Y$, which will be defined later in detail.

   Particularly, if $X$ and $Y$ are independent, $Cov(X, Y) = 0$ and
   \[
      Var(X + Y) = Var(X) + Var(Y)
   \]
\end{itemize}

\section{Independence}

Two events are independent if the outcome of one doesn't depend on the outcome of other.

For two independent events $A$ and $B$,

\[
   P(A \cap B) = P(A) \cdot P(B)
\]

Two random variables are independent if every event about $X$ is independent to every event about $Y$.

\[
   p(x, y) = p_X(x) \cdot p_Y(y)
\]

If two random variables $X$ and $Y$ are independent, we have the following properties:

\begin{align*}
   E[XY] &= E[X] \cdot E[Y] \\
   Var(X+Y) &= Var(X) + Var(Y) \\
   Cov(X, Y) &= 0
\end{align*}

\section{Conditional Probability and Bayes Theorem}

The conditional probability of $A$ given $B$ has occured (given $P(B) > 0$) is:

\[
   P(A | B) = \frac{P(A \cap B)}{P(B)}
\]

Also,

\[
   P(B | A) = \frac{P(A \cap B)}{P(A)}
\]

Rearranging, we have:

\[
   P(B) \cdot P(A | B) = P(A) \cdot P(B | A)
\]

Thus,

\[
   P(A | B) = \frac{P(A) \cdot P(B|A)}{P(B)}
\]

Above equation is known as the Bayes Theorem.

\section{Conditional Expectation}

\begin{itemize}
   \item The value $E[X | Y=y]$ is the expected value of $X$, given the value of $Y$ is known to be $y$.
   \item Thus, $E[X|Y]$ is a random variable itself, which is the expected value of $X$, choosing a random value for $Y$.
\end{itemize}

\subsection{Properties of Conditional Expectation}
\begin{itemize}
   \item{\bf Tower Property}
   \[
      E[E[X | Y]] = E[X]
   \]

   \begin{proof}
      \begin{align*}
         E[X | Y = y] &= \sum_x x \cdot P(X = x \, | \, Y = y) \\
         E[E[X | Y]] &= \sum_y E[X | Y = y] \cdot P(Y = y) \\
         &= \sum_y \left (\sum_x x \cdot P(X = x \, | \, Y = y)\right ) \cdot P(Y = y) \\
         \text{Now,} \\
         P(X = x | Y = y) &= \frac {P(X = x, Y = y)} {P(Y = y)} \\
         P(X = x, Y = y) &= P(X = x | Y = y) \cdot P(Y = y) \\
         \text{So,} \\
         E[E[X | Y]] &= \sum_y \left (\sum_x x \cdot P(X = x \, | \, Y = y)\right ) \cdot P(Y = y) \\
         &= \sum_y \left (\sum_x x \cdot P(X = x, Y = y) \right ) \\
         &= \sum_x \left (x \sum_y P(X = x, Y = y) \right ) \\
         &= \sum_x x \cdot P(X = x) \\
         &= E[X] \\
         E[E[X | Y]] &= E[X]
      \end{align*}
   \end{proof}

   \item {\bf Pull Out Property}
   \[
      E[h(Y) \cdot X | Y] = h(Y) \cdot E[X | Y]
   \]

   \begin{proof}
      \begin{align*}
         E[h(Y) \cdot X | Y = y] &= \sum_x h(y) \cdot x \cdot P(X = x \, | \, Y = y) \\
         &= h(y) \cdot \sum_x x \cdot P(X = x \, | \, Y = y) \\
         E[h(Y) \cdot X | Y = y] &= h(y) \cdot E[X | Y = y]
      \end{align*}
   \end{proof}

   \item {\bf Independence}
   If $X$ and $Y$ are independent,
   \[
      E[X | Y] = E[X]
   \]

   \begin{proof}
      \begin{align*}
         E[X | Y = y] &= \sum_x x \cdot P(X = x \, | \, Y = y) \\
         &= \sum_x x \cdot P(X = x) \\
         E[X | Y = y] &= E[X]
      \end{align*}
   \end{proof}
\end{itemize}

\end{document}





